\documentclass[main.tex]{subfiles}


\begin{document}

%from pagenumber to up
% \phantomsection 
% \hypertarget{toc}{}

 

 

 
   

\section{Конденсатор}


% \sbox{0}{%
% \includegraphics[]{el/15cap_pic.pdf}%
% }

% {%
% \setlength\intextsep{0pt}
% \begin{wrapfigure}[]{r}{\wd0}
%     \centering
%     \usebox{0}%
%     \caption{Плоский конденсатор}
%     \label{fig:p115}
% \end{wrapfigure}


\textbf{Конденсатор}~--- это устройство для накопления электрических зарядов. На рис.\,\ref{fig:p115} показана конструкция \emph{плоского конденсатора}.


    \begin{figure}[H]
    \centering
\includegraphics[]{el/15cap_pic.pdf}%
\caption{Плоский конденсатор}
\label{fig:p115}
\end{figure}


Простейший плоский конденсатор состоит из двух параллельных металлических пластин (\emph{обкладки}) площадью~$S$ каждая, расположенных на малом расстоянии~$d$ друг от друга и разделенных слоем диэлектрика\footnote{Также между обкладками может быть \emph{вакуум} ($\varepsilon=1$). Диэлектрическая проницаемость \emph{воздуха} обычно также считается равной единице.} с диэлектрической проницаемостью~$\varepsilon$. 
Конденсатор также имеет \emph{выводы} (черные линии на рис.\,\ref{fig:p115}) для его соединения с другими устройствами. 

Обычно в заряженном конденсаторе (то есть накопившем заряды на своих обкладках) одна из пластин конденсатора несет заряд~$+q$, а другая~--- заряд~$-q$. Под \emph{зарядом конденсатора} понимают заряд его \emph{положительной} обкладки.

\textbf{Емкость} ($C$ [Ф])~--- это характеристика конденсатора, показывающая его способность накапливать заряд:
\begin{equation}
    C=\frac{q}{U},
\end{equation}
где $q$~--- заряд конденсатора,
% \footnote{Обычно одна из пластин конденсатора несет заряд~$+q$, а другая~--- заряд $-q$.},
$U$~--- напряжение между пластинами.

Емкость численно равна заряду конденсатора при напряжении 1~В между его обкладками. 

\textbf{Емкость конденсатора} зависит от его размеров и диэлектрика в нем:
\begin{equation}
    C=\frac{\varepsilon\varepsilon_0S}{d},
\end{equation}
где $\varepsilon_0$~--- электрическая постоянная (см.~справочные таблицы).

\textbf{Энергия конденсатора} рассчитывается по формуле:
\begin{equation}
    W_C=\frac{CU^2}{2}.
\end{equation}

% \medskip

\noindent\textbf{Задача.} Плоский конденсатор состоит из двух параллельно расположенных в воздухе пластинок, каждая площадью 100~см$^2$, расстояние между ними 0,2~см. Определите емкость конденсатора.

\smallskip

\noindent\emph{Решение}. Между пластинами воздух: $\varepsilon=1$. Электрическая постоянная $\varepsilon_0\hm\approx8{,}85\cdot10^{-12}$~Ф/м. Тогда:
\begin{equation*}
    C=\frac{\varepsilon\varepsilon_0S}{d}=\frac{1\cdot8{,}85\cdot10^{-12}\cdot100\cdot10^{-4}}{0{,}2\cdot10^{-2}}=4{,}425\cdot10^{-11}~\text{Ф}.
\end{equation*}














 
\end{document}